\documentclass[11pt,a4paper]{article}
\usepackage[T1]{fontenc}
\usepackage[utf8]{inputenc}
\usepackage{amsmath,amssymb,amsthm}
\usepackage[margin=1in]{geometry}
\usepackage[colorlinks=true,linkcolor=blue,urlcolor=blue]{hyperref}
\theoremstyle{plain}
\newtheorem{theorem}{Theorem}
\newtheorem{lemma}[theorem]{Lemma}
\newtheorem{proposition}[theorem]{Proposition}
\newtheorem{corollary}[theorem]{Corollary}
\theoremstyle{definition}
\newtheorem{remark}[theorem]{Remark}

\title{The empirical recurrence for OEIS A269644, proved from an exact model}
\author{Adrian Perez Fontelles\\ \small Independent researcher}
\date{13 September 2026}
\begin{document}
\maketitle

\begin{abstract}
OEIS A269644 carries an empirical recurrence of order $8$. It is true, and it is
decidable rather than empirical. The entry's sequence is given exactly by an exact polynomial in $n$ past a computed transient, from
which $a$ provably satisfies a MONIC linear recurrence of order $S=37$, derived below and not
assumed. The residual of the conjectured recurrence satisfies the same one, so whether it
annihilates $a$ is settled by a finite exact computation. No transfer matrix is involved and
the count is not a walk.
\end{abstract}

\noindent\small 2020 Mathematics Subject Classification. 05A15, 05A19, 11P21.\normalsize

\section{The conjecture}

OEIS A269644 is ``Number of length-7 0..n arrays with no repeated value differing from the previous repeated value by other than plus two or minus 1.''. It has offset $1$ and begins
\[
20, 969, 9928, 55225, 216528, 675151, 1789528, 4200933,\ \dots
\]
The entry states, as a conjecture contributed by its author and never marked settled:
\begin{quote}
Empirical: a(n) = n\^{}7 + 7*n\^{}6 + 6*n\^{}5 + 20*n\^{}4 - 22*n\^{}3 + 28*n\^{}2 - 37*n + 13 for n>2. a(n) = 8*a(n-1) - 28*a(n-2) + 56*a(n-3) - 70*a(n-4) + 56*a(n-5) - 28*a(n-6) + 8*a(n-7) - a(n-8) for n>10.
\end{quote}
As of the ``Last modified'' line on the live entry (Jan 25 2019 11:44:42, revision 8)
this is still recorded as empirical, and nothing on the entry records it as proved.

\section{The count is a polynomial past a computed transient}

A REPEATED VALUE is a term equal to the one before it, and its value is that term; the entry's
condition relates each repeated value to the previous one. The length is fixed at $7$ and the
alphabet $0..n$ grows, so nothing here is a walk on a finite digraph --- the state that
remembers the previous repeated value has as many values as the alphabet. The count is taken
exactly instead. Writing $A[p][v]$ for the number of prefixes ending in $v$ whose previous
repeated value is $p$, appending a term $t$ gives
\[
A'[p][t] \mathrel{+}= \Big(\sum_v A[p][v]\Big) - A[p][t],
\qquad
A'[t][t] \mathrel{+}= A[p][t]\ \text{when the condition allows},
\]
the first for a $t$ differing from the last term and the second for a $t$ equal to it. Every
term the model produces is an exact integer.

\section{The bound}

Let $k$ be the point past which the condition stops caring: the predicate depends on the
difference of the two repeated values only through its sign and its magnitude up to $k$, and is
constant beyond that in each direction. Here $k=2$, read off the predicate rather than
assumed. An array is determined by its weak ordering --- the ordered set partition of the
$7$ positions into $m$ blocks of equal value, listed in increasing value --- together with
the base value $g_0\ge0$ and the gaps $g_1,\dots,g_{m-1}\ge1$ between consecutive distinct
values, subject to $g_0+\sum g_i\le n$. Each constraint is on a difference of two values, which
is a signed sum of consecutive gaps; since the predicate is constant past $k$, each constraint
splits into finitely many cases, and in every case it either fixes that gap-sum to one of at
most $2k+1$ values or pushes it beyond $k$. Within a case the count is the number of integer
points of a system (fixed amounts) plus (free gaps) $\le n$, a polynomial in $n$ of degree at
most $7$, valid as soon as $n$ exceeds the total forced amount. That total is at most
$(m-1)+(\text{constrained pairs})(k+1)\le L+L(k+1)$, so with
\[
n_0 \;=\; L + L(k+1) \;=\; 28
\]
the sequence agrees with a polynomial of degree at most $7$ for every $n\ge n_0$. It is
therefore annihilated by $(z-1)^{7+1}$ from that point, and by $z^{n_0+1}(z-1)^{7+1}$
everywhere, a monic recurrence of order $S=37$. The annihilator was then asked for terms the
model had not supplied and reproduced them, which is what a bound one short fails.

\section{The proof}

\begin{lemma}\label{lem:crit}
Let $A(z)=z^S-\sum_{j<S}\alpha_jz^{j}$ be the monic annihilator derived above and
$q(z)=z^{8}-\sum_ic_iz^{8-i}$ the characteristic polynomial of the conjectured
recurrence. The residual $r(n)=a(n)-\sum_ic_ia(n-i)$ is a linear combination of shifts of $a$,
so $A$ annihilates $r$ as well. Since $A(0)\ne0$ the recurrence it gives runs in both
directions, and $S$ consecutive zeros of $r$ force $r$ to vanish at every later index.
\end{lemma}

\begin{proof}
$A$ annihilates $a$, hence every shift of $a$, hence any linear combination of shifts; that is
$r$. A monic recurrence with nonzero constant term determines each term from the $S$ before it
and each from the $S$ after it, so a run of $S$ zeros propagates.
\end{proof}

\begin{theorem}
\[
a(n)\;=\;8\,a(n-1) - 28\,a(n-2) + 56\,a(n-3) - 70\,a(n-4) + 56\,a(n-5) - 28\,a(n-6) + 8\,a(n-7) - a(n-8)
\]
for every $n>10$.
\end{theorem}

\begin{proof}
The terms $a(0),\dots$ were computed exactly in integer arithmetic from the model, extended by
$A$ where the model itself was not evaluated, and $r(n)$ formed for each index. Those integers
vanish from $n=10+1$ onwards, and the run exceeds $S+8$, which by
Lemma~\ref{lem:crit} settles every larger index at once. The bound is exact: at $n=10$ the recurrence fails on the entry's own published terms, so no larger range holds.
\end{proof}

\section{Verification}

Three checks, each independent of the algebra above.

First, the model reproduces all $23$ terms the entry publishes, exactly, in integer
arithmetic. This is what ties the model to the entry: it is built by reading the entry's
English, and a misreading gives different counts.

Second, the conjectured recurrence was evaluated directly on those published terms, with no
model involved, and holds wherever the proved range and the data overlap.

Third, the derived annihilator $A$ was itself tested on terms it had not been given: the model
was evaluated beyond the $S$ values that determine it and $A$ reproduced them. A bound that is
too small fails this test, which is why it is run.

\begin{thebibliography}{9}
\bibitem{oeis} The OEIS Foundation, \emph{The On-Line Encyclopedia of Integer
Sequences}, \url{https://oeis.org}, 2026. Sequence A269644.
\bibitem{beckrobins} M.~Beck and S.~Robins, \emph{Computing the Continuous Discretely},
2nd ed., Springer, 2015. (Ehrhart quasi-polynomials and their periods.)
\bibitem{stanley} R.~P.~Stanley, \emph{Enumerative Combinatorics}, Volume 1, 2nd ed.,
Cambridge University Press, 2011.
\bibitem{ds} F.~Diamond and J.~Shurman, \emph{A First Course in Modular Forms},
Springer GTM 228, 2005.
\end{thebibliography}

\end{document}
